\chapter{四足运动控制导论与控制架构}
\label{ch:quad_arch}

% ════════════════════════════════════════════════════════════════
%  导论/架构章：四足机器人运动控制部 —— 全景分层架构与控制哲学。
%  主线 = 个人笔记《足式机器人控制》§概述/发展史/运动控制框架（控制=优化/规划问题、
%  连续单模型 vs 四足多模型须切换）；正史骨架 = 于宪元硕论 §1 绪论 + §2.1 FSM
%  （6 里程碑、浮动基欠驱、三问助记「在哪儿/去哪儿/怎么去」、FSM 阻尼保护）；
%  平台/组成 = 宇树《四足机器人控制算法》§1（典型平台谱、四系统组成）。
%  本章为「地图」：不做全量动力学推导（留各专章），公式极少且全为定位性/原理性，
%  硬件/电机/ROS 一律凝练进 practice。教学层遵循 v5.0 规范。
% ════════════════════════════════════════════════════════════════

\begin{practice}[前置自测]
开始之前，先试答下列问题。若已能流畅作答，可只速读本章架构图与速查表；若卡住，正是本章要为你解决的。
\begin{enumerate}
  \item 一辆自动驾驶汽车在高速路上的运动模型，与一只四足机器人从“站立”切到“小跑”再切到“跳跃”的运动模型，在“能否用\emph{一套不变的}动力学方程与约束描述”这一点上有什么本质区别？这个区别会逼出什么样的软件结构？
  \item 把“规划一条最省能的步态轨迹”说成“一个优化问题”，需要哪三样东西（目标、约束、变量）？当这个优化问题非线性、难以实时精确求解时，工程上通常做什么妥协？
  \item 四足机器人有 $12$ 个关节电机，但它的躯干在空间中有 $6$ 个自由度（位置 $3$ + 姿态 $3$）却\emph{没有}任何电机直接驱动。那么，控制器究竟靠“拧哪个螺丝”来让躯干达到期望的姿态与速度？
  \item 除了重力，一个在地面上行走的四足机器人，其躯干（浮动基座）受到的外力还有什么？这个事实为什么把控制范式从“控位置”推向了“控力”？
  \item 现代主流四足控制把“求多大的足底力”与“把这些力变成多大的关节力矩”分成了两层（MPC 与 WBC），而且让它们跑在\emph{不同}的频率上（如 $100\,\mathrm{Hz}$ 与 $500\,\mathrm{Hz}$）。为什么不合成一层、一个频率？
  \item 当机器人翻滚速度过快、或一条腿的跟踪误差大到危险时，控制系统应当“继续按原计划执行”还是“切换到某种安全行为”？这种“按情况切换行为”的逻辑，用什么计算结构来组织最自然？
\end{enumerate}
\end{practice}

\section*{本章目标}
学完本章，你应当能够：
\begin{enumerate}
  \item \textbf{说清}四足运动控制为何在本质上是一个\emph{带约束的优化/规划问题}，以及它与车辆/无人机的“连续单模型”相比，为何是一个必须\emph{切换多模型}的问题（\cref{sec:qa-why-hard}）；
  \item \textbf{论证}为什么“固定连杆机构”“纯位置控制”“纯瞬时（无预测）控制”这三种朴素思路都不足以驱动动态四足，并由此\emph{逼出}力控、预测与分层（\cref{sec:qa-naive}）；
  \item \textbf{复述}足式控制理论的六个历史里程碑（固定连杆 $\to$ 位置控制 $\to$ 人类驾驶 $\to$ Raibert 三通道/SLIP $\to$ Pratt VMC $\to$ WBC/MPC），并在典型平台谱（波士顿动力、MIT、ETH、IIT、UPenn、宇树）上定位（\cref{sec:qa-history}）；
  \item \textbf{画出并解释}四足动态运动控制的\emph{分层数据流}：状态估计 $\to$ 质心/足底规划 $\to$ MPC（求虚拟足底力）$\to$ WBC（转关节力矩）$\to$ 电机，并用“在哪儿/去哪儿/怎么去”三问作助记主轴（\cref{sec:qa-arch}）；
  \item \textbf{论证}浮动基座的\emph{六自由度欠驱}本质，说明“地面反力是除重力外唯一外力”如何导出\emph{力控范式}，并用虚功一句话定位“足底力 $\to$ 关节力矩”的最后一公里 $\boldsymbol\tau=\mathbf{J}^\top\mathbf{f}$（\cref{sec:qa-arch}）；
  \item \textbf{用有限状态机（FSM）}组织起立/下蹲/行走/平衡站立等行为的合法切换，理解“忙碌时间”“危险检测 $\to$ 阻尼保护”这套安全机制，并写出阻尼保护控制律 $\boldsymbol\tau_{\mathrm{motor}}=-k_{\mathrm{damping}}\,v_{\mathrm{motor}}$（\cref{sec:qa-fsm}）；
  \item \textbf{辨析}“优化”与“强化学习”两大流派的本质共性（都是在某函数下求最优运动策略）与分工，理解 MPC 的“接触-规划分离”为何降复杂度、又为何难以扩展到接触不确定性（\cref{sec:qa-two-streams}）；
  \item \textbf{在架构图上为后续每一章“点名定位”}——运动学、步态、状态估计、单刚体 MPC、WBC、实践各落在哪一层（\cref{sec:qa-arch,sec:qa-map}）。
\end{enumerate}

\subsection*{本章知识导航}
本章是整部“四足运动控制”的\emph{地图}。它不推导任何一条完整的动力学或优化方程（那些是各专章的任务），而是回答一个统领性的问题：\textbf{为什么动态四足的控制系统会长成“估计 $\to$ 规划 $\to$ MPC $\to$ WBC $\to$ 电机”这样一条分层流水线，外加一个切换行为的状态机？}结构如下：

\begin{center}
\begin{forest} navtreeLR
[{四足运动控制\\导论与架构}
  [{为何难且特别\\优化问题+多模型}
    [{连续单模型\\vs 多模型切换}] ]
  [{朴素思路为何不行\\连杆/位置/瞬时}
    [{逼出力控+预测+分层}] ]
  [{历史与平台谱\\6 里程碑}
    [{Raibert/VMC/WBC/MPC}] ]
  [{分层架构\\三问助记}
    [{估计$\to$规划$\to$MPC$\to$WBC}] ]
  [{相位组织\\FSM+阻尼保护}
    [{优化 vs RL 两流派}] ]
]
\end{forest}
\end{center}

\noindent\textbf{推荐路径}：\cref{sec:qa-why-hard,sec:qa-naive} 建立“控制即优化、四足须切换多模型”和“朴素思路为何失败”的两大立论，是全章的哲学地基；\cref{sec:qa-history} 给出从连杆机构到 MPC$+$WBC 的正史脉络与平台实物锚，可作背景速读；\cref{sec:qa-arch} 是\emph{本章核心}——把分层数据流逐层讲清并为每层交叉引用到专章；\cref{sec:qa-fsm} 用 FSM 组织相位切换、落地“多模型须切换”这一主线；\cref{sec:qa-two-streams} 点出“优化 vs 强化学习”的现代分叉作收束；\cref{sec:qa-map} 是一张“后续章节在架构图上的位置”总表，可随时回查。

\subsection*{前置知识桥接}
本章是导论，几乎不依赖具体数学工具，但它\emph{指向}的每一层都对应全书前几部的地基：
\begin{itemize}
  \item \cref{ch:control}（控制导论）给出“代价函数 + 约束 + 滚动优化”的一般框架，本章把四足控制\emph{装进}这一框架后再分层；
  \item \cref{ch:lie}（李群与李代数）、\cref{ch:rigid_body}（刚体运动）给出 $\SO(3)$、$\SE(3)$ 与刚体位姿的语言，本章谈“浮动基座”“躯干姿态”时直接借用；
  \item \cref{ch:nlopt}（非线性优化/QP）是 MPC 与 WBC 落地为可解问题的求解器基础，本章只“点名”不展开；
  \item \cref{ch:rl-basic}（强化学习导论）是本章末“另一条腿”——优化 vs RL 两流派的 RL 一侧的入口。
\end{itemize}

\subsection*{如果跳过本章会怎样}
\begin{itemize}
  \item 你会在后续各专章里学到 MPC、WBC、状态估计、步态规划的\emph{局部}细节，却看不清它们如何\emph{拼成一台会跑会跳的机器人}——每一章都像孤岛，不知道上游给它什么、下游要它给什么；
  \item 你会理解“四足为什么能动”停留在“迈腿 + 平衡”的直觉，无法解释为什么现代系统\emph{必须}分层、\emph{必须}力控、\emph{必须}有状态机——这正是本章 \cref{sec:qa-naive,sec:qa-arch,sec:qa-fsm} 的立论。
\end{itemize}

\begin{table}[H]\centering\small
\caption{本章阅读方式建议}
\begin{tabular}{lll}
\toprule
阅读方式 & 时间 & 适合谁 \\
\midrule
精读（含历史与全部洞见） & 2--3 小时 & 首次进入四足运动控制、要建立全局观 \\
速读（跳过历史与平台谱） & 1 小时 & 有控制基础、只想对齐本书架构与记号 \\
速查（只看架构图 + 章节定位总表） & 10 分钟 & 读专章时回来查“这一层在哪、接口是什么” \\
\bottomrule
\end{tabular}
\end{table}

\begin{note}
\textbf{本章记号与约定.}\;
本章为导论，公式极少。沿用全书约定：旋转矩阵 $\mathbf{R}\in\SO(3)$、本书几何主线取\emph{右扰动}、四元数取 Hamilton 约定、$\se(3)$ 平移分量在前（\cref{ch:lie,ch:rigid_body}）。
机身（浮动基）记 $\{b\}$、世界（惯性）系记 $\{s\}$ 或 $\mathcal{W}$；机身到世界的旋转记 $\mathbf{R}_{sb}$。
第 $i$ 腿足底反力（GRF）记 $\mathbf{f}_i$，关节力矩记 $\boldsymbol{\tau}$，单腿雅可比记 $\mathbf{J}_i$；单腿三关节统一记 abad/hip/knee（与广义坐标 $\mathbf q$ 的腿关节分量 $q_i$ 一致，\cref{ch:quad_kin}）。
控制频率：MPC 周期 $\Delta T$、其余控制环周期 $\Delta t$（沿用 \cref{ch:quad_mpc} 的记号，与足端力 $f$ 不撞名）。
反对称化算子统一记 $(\cdot)^\wedge$、其逆 $(\cdot)^\vee$（\cref{ch:lie}）。
\end{note}

% ════════════════════════════════════════════════════════════════
\section{动机：四足控制为什么“难且特别”}
\label{sec:qa-why-hard}

\paragraph{第一性原理：控制即优化/规划。}
要把整部书放进一个统一的视角，先从最高处看一句话：\textbf{足式运动控制本质上是一个带约束的优化/规划问题}。给定一个性能指标（路径最短、速度最快、能耗最小、跟踪误差最小等），在一组约束（加速度/速度上限、关节限位、摩擦不打滑、不可涉足的空间等）下，求出令指标最优的\emph{控制量}（在四足里，最终落到关节力矩，中间常表达为足底反力或落足点）。这就是把控制问题“写成数学”的统一姿态：目标函数 + 约束 + 决策变量\cite{yu2021quadruped}。

\begin{insight}[控制 $=$ 优化/规划，难解时牺牲解空间换可行解（全章总纲）]\label{ins:qa-control-is-opt}
全书控制哲学的开篇总纲是一句话：\textbf{足式运动控制本质是一个优化/规划问题}——在“最短/最快/最省能”等指标下、满足“加速度与速度上限、不可涉足空间”等约束，求最优控制量。但现实里多数情况\emph{无法对系统完美建模}，或问题\emph{非线性而难以精确实时求解}；于是工程上做一个反复出现的妥协：\textbf{牺牲一部分解空间（用简化模型、凸近似、有限时域、固定步态时序等手段缩小搜索范围），换取一个能在控制周期内算出来的\emph{次优但可行}的解}。这条“宁要可行的次优、不要无法实时的最优”的取舍，是理解后续每一章为何要做各种近似（单刚体、凸化、摩擦金字塔、有限预测时域）的总钥匙。
\end{insight}

\begin{insight}[连续单模型 vs 多模型须切换（导出 FSM）]\label{ins:qa-multimodel}
为什么四足比汽车、无人机“特别”？关键在一个对照：\textbf{车辆与无人机的运动是“连续单模型”}——它们的运动模式与约束在整个任务里几乎\emph{不变}，可以用\emph{一套}动力学方程与一套约束统一描述、统一控制。而\textbf{四足是“多模型须切换”}：它存在站立、支撑、摆动、飞行（腾空）等\emph{多种本质不同的运动模型}，每种模型的接触集、可控自由度、动力学结构都不一样，因此控制系统\emph{必须}在这些模型之间做\emph{显式的模式切换}。这一点直接导出本章 \cref{sec:qa-fsm} 的有限状态机：当“一套模型走天下”不再成立，就需要一个机构来管理“此刻该用哪套模型、何时允许切到下一套”。这是“四足控制为什么长成现在这样”的第一个分叉。
\end{insight}

\paragraph{足式相对轮/履的两大本质优势。}
既然四足更难，为何还要用腿？有两条\emph{本质}优势（不是工程细节，而是运动学/动力学层面的结构性好处）：
\begin{enumerate}
  \item \textbf{非连续支撑}：轮式/履带式需要\emph{连续}的可支撑表面，而腿式只需要一串\emph{离散}的落足点。这使腿式能跨越沟壑、台阶、碎石等非结构化地形——典型如爬楼梯：腿可以“点”在每一级台阶的边缘，而轮子需要连续斜坡\cite{yu2021quadruped}。
  \item \textbf{本体与腿运动解耦}：腿式系统的躯干运动与腿的运动在一定程度上\emph{解耦}，腿可以像\emph{主动悬挂}一样吸收地形起伏，让躯干（及其搭载的传感器/负载）保持平稳——这是轮式底盘难以做到的避振能力\cite{yu2021quadruped}。
\end{enumerate}
代价也很明确：腿式结构复杂（多关节、多电机）、能耗高、续航短\cite{yu2021quadruped}。这正是 \cref{ins:qa-control-is-opt} 里“能耗最小”常被列入指标的现实原因。

\begin{figure}[ht]
\centering
\begin{tikzpicture}[scale=0.95,>=Stealth]
  % --- left: wheeled needs continuous ramp ---
  \begin{scope}[xshift=-3.6cm]
    \node[font=\small\bfseries] at (0,2.2) {轮式：需连续支撑};
    % stairs as obstacle
    \draw[gray!50,fill=gray!12] (-1.6,-1.0) -- (0.2,-1.0) -- (0.2,-0.6) -- (0.7,-0.6) -- (0.7,-0.2) -- (1.2,-0.2) -- (1.2,0.2) -- (1.6,0.2) -- (1.6,-1.0) -- cycle;
    % wheel stuck at first step
    \draw[main!60!black, very thick, fill=main!10] (0.0,-0.7) circle (0.3);
    \draw[red!70!black,-{Stealth[length=2mm]},very thick] (0.35,-0.55) -- (0.6,-0.3);
    \node[red!70!black,font=\scriptsize] at (0.95,-0.05) {卡住};
  \end{scope}
  % --- right: legged steps on discrete points ---
  \begin{scope}[xshift=3.0cm]
    \node[font=\small\bfseries] at (0,2.2) {足式：点离散支撑点};
    \draw[gray!50,fill=gray!12] (-1.6,-1.0) -- (0.2,-1.0) -- (0.2,-0.6) -- (0.7,-0.6) -- (0.7,-0.2) -- (1.2,-0.2) -- (1.2,0.2) -- (1.6,0.2) -- (1.6,-1.0) -- cycle;
    % body
    \draw[rounded corners=2pt, fill=main!8, draw=main!60!black] (-0.5,0.9) rectangle (0.9,1.3);
    % two legs to discrete edges
    \draw[gray!60!black,thick] (-0.4,0.9) -- (-0.9,-0.6);
    \fill (-0.9,-0.6) circle (2pt);
    \draw[gray!60!black,thick] (0.8,0.9) -- (1.0,-0.2);
    \fill (1.0,-0.2) circle (2pt);
    \node[second,font=\scriptsize] at (-0.9,-0.95) {落足点};
    \node[second,font=\scriptsize] at (1.05,0.1) {落足点};
  \end{scope}
  \draw[gray] (-5.3,-1.0)--(4.7,-1.0);
\end{tikzpicture}
\caption{足式相对轮式的“非连续支撑”优势：轮式需要连续可支撑表面，遇台阶易卡住；足式只需把脚“点”在离散落足点上即可跨越非结构化地形。这是腿式不可替代性的运动学根源。}
\label{fig:qa-legged-advantage}
\end{figure}

\paragraph{解的去向：估计 $\to$ 规划 $\to$ 控制的闭环。}
把“控制即优化”落到运行时，还需看清解的\emph{流向}：上层优化/规划器求出的“解”通常是一条\emph{期望轨迹}（质心轨迹、足端轨迹），它被\emph{发给底层的跟踪控制器}去执行；与此同时，一个\emph{观测器/状态估计器}（融合 IMU、轮速计/足底里程计等）\emph{实时}估计机器人当前的位置、速度、加速度，反馈回规划与控制\cite{yu2021quadruped}。这就形成了贯穿全书的闭环数据流：\textbf{估计当前在哪儿 $\to$ 规划要去哪儿 $\to$ 控制怎么去}——它正是 \cref{sec:qa-arch} 分层架构的雏形，也对应“三问”助记（\cref{ins:qa-three-questions}）。

\begin{practice}
\begin{enumerate}
  \item 用 \cref{ins:qa-control-is-opt} 的“目标 + 约束 + 变量”三件套，把“让四足以 $1\,\mathrm{m/s}$ 匀速直行且尽量省电”写成一个优化问题的\emph{文字版}（不必写出方程）。指出哪一项是目标、哪些是约束。
  \item 举一个生活中的“连续单模型”系统与一个“多模型须切换”系统（不必是机器人），并用 \cref{ins:qa-multimodel} 的语言说明区别。
  \item （概念）若把四足强行用“一套不变的动力学模型 + 一套约束”来控制（无视支撑/摆动/飞行的差异），最可能在什么相位、什么时刻出问题？
\end{enumerate}
\end{practice}

% ════════════════════════════════════════════════════════════════
\section{反面：朴素思路为什么不行}
\label{sec:qa-naive}

理解一个架构为何\emph{这样}设计，最快的路径是看清\emph{别的设计为何失败}。本节按历史顺序排出四种朴素思路，每一种的失败都\emph{逼出}现代架构的一个支柱：力控、预测、分层。

\paragraph{反面一：固定运动模式的连杆机构。}
最早的“行走机器”（如基于切比雪夫连杆等机构的步行装置）用\emph{固定}的连杆几何把电机的旋转变成腿的周期摆动。其根本缺陷是：腿的轨迹被机构\emph{写死}，机身一旦受扰或地形变化，系统\emph{无法响应}、不能自适应。早在 $20$ 世纪 $50$ 年代人们就认识到：\textbf{真正的足式机器人必须有反馈系统}\cite{yu2021quadruped}。

\begin{pitfall}[历史教训：无反馈的固定机构不是“机器人”]
固定连杆机构能“走”，但它\emph{不能控}——它没有传感、没有反馈、没有对扰动和地形的响应能力。把“能周期性迈腿”误当成“会行走的机器人”，是足式发展史最早的弯路。现代四足的每一层都建立在“感知—反馈—响应”之上，这正是它区别于一台机械步行玩具的根本。
\end{pitfall}

\paragraph{反面二：纯位置控制。}
一个看似自然的升级是：规划好足底轨迹 $\to$ 逆运动学解出关节角 $\to$ 把关节角\emph{位置}指令发给电机的位置环。但这与 $19$ 世纪的连杆机器人\emph{没有本质区别}：电机内部的位置环是一个\emph{黑箱}，当腿与环境发生\emph{力的作用}（触地冲击、地形高低差、外推）时，纯位置控制的行为\emph{不可控}——它会“硬顶”地面或在悬空时“踏空”。更糟的是，位置环需要\emph{响应时间}才能消除误差，这\emph{压低了控制带宽}，难以应对高动态运动\cite{yu2021quadruped}。

\begin{pitfall}[误区：位置控制的“黑箱力”]
位置控制只承诺“关节角到达期望值”，却\emph{不承诺}过程中的\emph{力}。一旦腿与地面有力交互，位置环为了“到位”可能输出任意大的力（撞地）或在落空时输出零力（踏空）——这两者都不可控。结论：与环境有力交互的腿式系统，控制的\emph{第一公民应当是力而非位置}。这把范式从“控位置”推向了“控力”（\cref{ins:qa-floating-base,ins:qa-grf-only}）。
\end{pitfall}

\paragraph{反面三：只用瞬时/无预测控制器。}
即便用上了力控（例如纯 WBC 或瞬时 QP，在每个控制周期把当前任务映射成当前关节指令），若控制器\emph{没有预测能力}，在动态步态（尤其含腾空相）下仍会失败：在\emph{支撑相的后半段}，无预测控制器\emph{不会为即将到来、没有地面反力的飞行相做准备}；在\emph{飞行相的后半程}，它又\emph{不会为后续的落地冲击做准备}\cite{yu2021quadruped}。这条短板在 \cref{ch:quad_mpc} 的“间歇欠驱”立论里被推到极致，并\emph{逼出}了对模型预测控制（MPC）的需求。

\begin{insight}[无预测则相位切换处崩溃，故需 MPC（串 \cref{ch:quad_mpc}）]\label{ins:qa-need-mpc}
\textbf{瞬时控制器（PID、纯 WBC、瞬时 QP）是反应式的——误差出现才纠正、只看当前，无法为未来的相位切换“排兵布阵”}。而四足动态步态的相位\emph{是已知的}（步态调度器预先给出谁何时触地、何时腾空），这份“已知的未来”恰好可以被 MPC 的预测时域利用起来：在可驱的窗口里为即将到来的欠驱窗口（飞行相、对角支撑的欠驱方向）攒够初速度与姿态。这就是“为什么架构里要塞一个 MPC 慢环”的根因。完整的间歇欠驱论证、单刚体建模与凸 QP 推导见 \cref{ch:quad_mpc}；本章只把它\emph{钉}在架构图的“怎么去”那一层。
\end{insight}

\paragraph{反面四：参数过少、手工设计的早期学习方法。}
作为“另一条腿”的预告：早期的学习/强化学习方法因\emph{模型参数少、特征靠手工设计、需在硬件上反复微调}，表征能力受限，难以直接产出复杂动态步态\cite{yu2021quadruped}。这一短板要等到深度网络 + GPU 并行模拟器（见 \cref{sec:qa-two-streams}）才被突破。它解释了为何很长一段时间里，四足动态控制的主流是\emph{优化/MPC}而非学习。

\begin{practice}
\begin{enumerate}
  \item 把四个“反面”各对应到它\emph{逼出}的一个现代架构支柱（力控 / 预测 / 分层 / 学习），列成一张两列表。
  \item （概念）一个只跟踪机身高度的纯 PID 四足，在从“小跑”切到“含腾空的飞奔”时，最可能在落地瞬间发生什么？用 \cref{ins:qa-need-mpc} 的语言解释。
  \item 解释“位置环响应时间压低控制带宽”这句话：为什么“需要时间消除误差”会限制系统能跟踪的最高频率的运动？
\end{enumerate}
\end{practice}

% ════════════════════════════════════════════════════════════════
\section{历史与平台谱：从连杆到 MPC\texorpdfstring{$+$}{+}WBC}
\label{sec:qa-history}

把上一节的四个反面接起来，就是足式控制理论的\emph{正史}。沿六个里程碑逐一定位，并配以典型平台谱与三线（硬件/模拟器/算法）演进，让“架构为何如此”有历史纵深\cite{yu2021quadruped}。

\subsection{六个里程碑}
\label{subsec:qa-milestones}

\begin{figure}[ht]
\centering
\begin{tikzpicture}[>=Stealth, font=\scriptsize]
  \draw[thick,-{Stealth[length=2.5mm]}] (-0.3,0) -- (12.3,0);
  \foreach \x/\lab in {0.5/{1.}, 2.5/{2.}, 4.5/{3.}, 6.5/{4.}, 8.5/{5.}, 10.5/{6.}} {
    \fill[main!70!black] (\x,0) circle (2pt);
  }
  \node[align=center,anchor=south] at (0.5,0.15) {\textbf{固定连杆}\\1870s\\无反馈};
  \node[align=center,anchor=north] at (2.5,-0.15) {\textbf{位置控制}\\足底轨迹\\仍无力控};
  \node[align=center,anchor=south] at (4.5,0.15) {\textbf{人类驾驶}\\GE 四足卡车\\闭环靠人};
  \node[align=center,anchor=north] at (6.5,-0.15) {\textbf{Raibert 三通道}\\1980s\,/\,SLIP\\精准动力学非必要};
  \node[align=center,anchor=south] at (8.5,0.15) {\textbf{Pratt VMC}\\2001\\虚拟组件};
  \node[align=center,anchor=north] at (10.5,-0.15) {\textbf{WBC\,/\,MPC}\\优先级+预测\\现代主流};
\end{tikzpicture}
\caption{足式控制理论的六个里程碑\cite{yu2021quadruped}：从无反馈的固定连杆机构，到有足底轨迹但无力控的位置控制，到靠人闭环的液压四足卡车，到 Raibert 用解耦三通道与 SLIP 模型证明“精准动力学非绝对必要”，到 Pratt 的虚拟模型控制，最终汇入今天的 WBC（优先级任务分解）+ MPC（预测控制）现代局面。}
\label{fig:qa-timeline}
\end{figure}

\paragraph{\cnum{1} 固定连杆机构。}
摄影师迈布里奇（Eadweard Muybridge）于 $1878$ 年用连续抓拍证明了马奔跑时\emph{四蹄同时离地}的腾空相，激发了对足式运动机制的研究；$19$ 世纪 $70$ 年代切比雪夫等以连杆机构实现周期步行的装置随之出现。但如 \cref{sec:qa-naive} 所述，它们\emph{无反馈}，只是“会迈腿的机构”\cite{yu2021quadruped}。

\paragraph{\cnum{2} 可自由设计足底轨迹的位置控制。}
比固定机构进了一步：腿的足底轨迹\emph{可由设计者自由规划}（不再被连杆几何写死），再用逆运动学转成关节位置指令。但其控制\emph{本质仍是位置控制}，缺乏力控（\cref{sec:qa-naive} 反面二），与环境有力交互时不可控\cite{yu2021quadruped}。

\paragraph{\cnum{3} 由训练有素的人类驾驶。}
$20$ 世纪中叶的 GE 四足卡车（液压驱动）通过\emph{力反馈把手}让人类驾驶员操纵每条腿——这是历史上第一次出现“力反馈 + 闭环”的足式系统，但\emph{闭环靠人}，驾驶员极易疲劳，无法自主\cite{yu2021quadruped}。

\paragraph{\cnum{4} Raibert 三通道解耦控制与 SLIP 模型。}
$20$ 世纪 $80$ 年代，Marc Raibert 的奠基工作\cite{raibert1986legged} 提出把控制\emph{解耦成三个独立通道}：\emph{跳跃高度}（由支撑相蹬地力调）、\emph{机身姿态}（由支撑相足底力矩调）、\emph{前进速度}（由\emph{飞行相落足点}调，落点取为机身速度的近似线性函数）。更深刻的是他把单足/双足/四足都抽象为\textbf{弹簧负载倒立摆（SLIP）}模型，用极简模型在多种地形上实现动态运动，从而论证了一个影响深远的结论：\textbf{精准的全身动力学并非动态运动的绝对必要条件}\cite{yu2021quadruped}。这条理念至今仍是凸 MPC 的落足初值与对照基线。

\begin{remark}[SLIP 概念模型的示意形式]\label{rem:qa-slip}
为帮助理解，这里以\emph{标准 SLIP 形式}给出其概念方程（点质量 $m$ + 无质量弹性腿）：
\begin{equation}
  m\,\ddot{\boldsymbol r}=\boldsymbol F_{\mathrm{leg}}+m\boldsymbol g,
  \qquad
  \boldsymbol F_{\mathrm{leg}}=k_{\mathrm{leg}}\,(\ell_0-\ell)\,\hat{\boldsymbol n},
  \label{eq:qa-slip}
\end{equation}
其中 $\boldsymbol r$ 为质心位置，$k_{\mathrm{leg}}$ 为腿刚度、$\ell_0$ 为自然腿长、$\ell$ 为当前腿长、$\hat{\boldsymbol n}$ 为沿腿方向的单位矢量。\cref{eq:qa-slip} 按 SLIP 标准定义\cite{raibert1986legged} 给出，仅供建立直觉，是概念示意而非完整动力学。SLIP 的完整动力学与四足简化模型谱系见 \cref{ch:quad_mpc} 的三模型对照。
\end{remark}

\paragraph{\cnum{5} Pratt 虚拟模型控制（VMC）。}
$2001$ 年，Pratt 等提出\textbf{虚拟模型控制}\cite{pratt2001virtual}：在机器人与环境之间\emph{想象}一组虚拟的弹簧、阻尼、质量等机械组件，这些虚拟组件产生的虚拟力/力矩再通过雅可比转置映射成\emph{等效关节力矩}。其优点是\emph{计算量小}、且可\emph{自由设计}虚拟组件的种类与参数来塑造期望的交互行为\cite{yu2021quadruped}。VMC 是“用直觉的力学元件设计力控”的代表，其思想被现代 WBC/MPC 部分吸收。

\paragraph{\cnum{6} WBC 与 MPC。}
现代两大支柱：
\begin{itemize}
  \item \textbf{全身控制 WBC}：把控制目标拆成多个\emph{任务}并赋予\emph{优先级}，用\emph{零空间投影}保证低优先级任务\emph{不干扰}高优先级任务；本书的 WBC 至少实现\emph{浮动基逆动力学}（把期望加速度/力转成关节力矩）\cite{yu2021quadruped}。但 WBC \emph{无预测性}（\cref{ins:qa-need-mpc}）。
  \item \textbf{模型预测控制 MPC}：用\emph{滚动时域优化}跟踪期望轨迹，\emph{弥补 WBC 的无预测短板}；代价是\emph{在线求解大规模优化、算力消耗高、控制带宽受限}\cite{yu2021quadruped}。
\end{itemize}

\begin{insight}[演进总趋势——从无到有、从简化到全动力学、从瞬时到预测]\label{ins:qa-evolution}
把六个里程碑拉直，可见一条清晰的演进主轴：\textbf{从无反馈到有反馈；从位置控制到力控；从简化模型（SLIP/VMC）到全（多刚体）动力学（WBC）；从瞬时控制到预测控制（MPC）}。今天的现代局面正是这条主轴的汇流：\textbf{以 MPC 求足底反力 + 以 WBC（含多刚体动力学）转关节力矩 + 借鉴 VMC/SLIP 的力控理念}\cite{yu2021quadruped}。这也正是本部各专章如何\emph{拼成一张图}的总声明——MPC 是 \cref{ch:quad_mpc}，WBC 是 \cref{ch:quad_wbc}，二者之上还有状态估计 \cref{ch:quad_est}、步态规划 \cref{ch:quad_gait}、运动学 \cref{ch:quad_kin}。
\end{insight}

\subsection{典型平台谱}
\label{subsec:qa-platforms}

历史的另一条线是\emph{实物平台}。从研究室到产品的平台谱可作“理论落在哪些机器上”的实物锚（凝练登记，细节见各平台公开资料）\cite{unitree2023quadruped}。

\begin{table}[H]\centering\small
\caption{典型四足平台谱（附控制史脉络）\cite{unitree2023quadruped}}
\begin{tabular}{p{0.20\textwidth} p{0.30\textwidth} p{0.40\textwidth}}
\toprule
机构 / 谱系 & 代表平台 & 控制史定位 \\
\midrule
波士顿动力 & BigDog、Atlas、Spot、Stretch & 源出 MIT Leg 实验室（Raibert），液压起家、动态平衡标杆 \\
MIT Biomimetic & Cheetah 3、Mini Cheetah & 单刚体凸 MPC 的硬件载体，首个实现后空翻的四足，软硬件开源 \\
ETH RSL & ANYmal & 高保真力控 + 学习/优化融合，工业巡检落地 \\
IIT DLS & HyQ & 液压驱动四足，强力交互与冲击适应 \\
UPenn GRASP / Ghost Robotics & Minitaur、Vision 系列 & 直驱/准直驱低成本动态四足 \\
宇树 Unitree & Laikago $\to$ Aliengo / A1 / Go1 & Laikago 为首台公开零售四足，推动低成本高性能普及 \\
\bottomrule
\end{tabular}
\label{tab:qa-platforms}
\end{table}

\paragraph{四系统组成（凝练）。}
一台四足机器人可粗分为四个系统：\emph{动力系统}（关节电机及传动）、\emph{控制系统}（运算单元 + 控制算法）、\emph{通信系统}（机内/机外数据链路）、\emph{能源系统}（电池及管理）\cite{unitree2023quadruped}。典型构型为 \textbf{$12$ 个关节}（每腿 $3$ 个：髋外摆 abad、髋俯仰 hip、膝 knee），由\emph{永磁同步电机}驱动，可做精确的力矩/角度控制——这正是 \cref{sec:qa-arch} 力控范式得以落地的\emph{硬件前提}。

\begin{insight}[硬件是地基——足式发展离不开硬件，如 AI 离不开芯片]\label{ins:qa-hardware}
足式机器人的发展离不开硬件的发展，正如 AI 离不开芯片。具体而言，\emph{现代高功率密度、可控扭矩的电机}（准直驱/低减速比 + 高带宽力矩控制）是低成本、高性能四足得以普及的关键使能技术——没有能“精确、快速地输出期望力矩”的电机，\cref{sec:qa-arch} 里“算出足底力 $\to$ 发关节力矩”的整条力控链都无从落地\cite{yu2021quadruped,unitree2023quadruped}。这也与 \cref{ch:quad_mpc} 提到的“硬件配合”（把驱动电机上移髋部、降低腿部转动惯量以逼近单刚体假设）形成呼应：\textbf{算法的简化假设，需要硬件设计来“补偿”才成立}。
\end{insight}

\begin{practice}
\begin{enumerate}
  \item 在 \cref{fig:qa-timeline} 的六个里程碑上，标出每一个分别引入了“反馈 / 力控 / 解耦 / 优先级 / 预测”中的哪一项能力。
  \item 选 \cref{tab:qa-platforms} 中任一平台，查其公开资料，说明它主要用“优化/MPC”还是“强化学习”路线（或两者），并对照 \cref{sec:qa-two-streams}。
  \item （概念）用 \cref{ins:qa-hardware} 解释：为什么“高功率密度可控扭矩电机”被称为四足普及的\emph{使能技术}，而不仅仅是一个零件升级？
\end{enumerate}
\end{practice}

% ════════════════════════════════════════════════════════════════
\section{理论锚定：四足动态运动控制的分层架构}
\label{sec:qa-arch}

这是本章的\emph{核心}。以“三问”作主轴，把运动控制框架展开为一条\emph{分层数据流}，并为每一层交叉引用到对应专章——这就是全部书的“地图”\cite{yu2021quadruped}。

\subsection{架构总览与三问助记}
\label{subsec:qa-overview}

\begin{insight}[三问——在哪儿 / 去哪儿 / 怎么去（分层主轴）]\label{ins:qa-three-questions}
一组朴素的设问可以把整个控制系统串起来，本书采为分层架构的\emph{助记主轴}：
\begin{enumerate}
  \item \textbf{“在哪儿？”} $\to$ \emph{状态估计}：机器人此刻的位姿、速度、角速度是多少？
  \item \textbf{“去哪儿？”} $\to$ \emph{规划}：期望的质心轨迹与各足端轨迹是什么？
  \item \textbf{“怎么去？”} $\to$ \emph{控制（MPC + WBC）}：施加多大的足底力、最终发多大的关节力矩，才能把当前状态驱往期望轨迹？
\end{enumerate}
这三问\emph{一一对应} \cref{sec:qa-why-hard} 末尾“估计 $\to$ 规划 $\to$ 控制”的闭环数据流，是理解后续每一章\emph{属于哪一问}的总索引\cite{yu2021quadruped}。
\end{insight}

\cref{fig:qa-pipeline} 给出完整的分层数据流。\textbf{实线}表示\emph{控制信号流}（指令逐层下传），\textbf{虚线}表示\emph{状态信号流}（估计结果反馈给各层）。最上游是外部输入：\emph{路径规划器/遥控器}给出期望速度指令，\emph{步态调度器}给出步态指令（谁何时支撑/摆动）；最下游是\emph{电机}\cite{yu2021quadruped}。

\begin{figure}[ht]
\centering
\begin{tikzpicture}[>=Stealth, node distance=6mm,
  blk/.style={draw, rounded corners, align=center, font=\small, inner sep=4pt, minimum height=0.95cm, minimum width=2.5cm},
  ctrl/.style={-{Stealth[length=2.5mm]}, thick, main!60!black},
  state/.style={-{Stealth[length=2.5mm]}, thick, second!70!black, dashed}]
  % top inputs
  \node[blk, fill=gray!8] (cmd) {上层指令\\（速度 + 步态）};
  % main vertical chain
  \node[blk, fill=second!10, draw=second!60!black, below=of cmd] (plan) {规划\\质心轨迹 + 足端轨迹\\\scriptsize“去哪儿”};
  \node[blk, fill=main!10, draw=main!60!black, below=of plan] (mpc) {MPC（慢环 $\Delta T$）\\单刚体 $\to$ 求足底反力 $\mathbf{f}_i$\\\scriptsize“怎么去·上”};
  \node[blk, fill=main!10, draw=main!60!black, below=of mpc] (wbc) {WBC（快环 $\Delta t$）\\多刚体逆动力学 $\to$ 关节力矩\\\scriptsize“怎么去·下”};
  \node[blk, fill=gray!8, below=of wbc] (motor) {电机\\（永磁同步、力矩控制）};
  % robot/plant
  \node[blk, fill=third!10, draw=third!60!black, below=of motor] (robot) {机器人本体 + 环境};
  % state estimator on the left
  \node[blk, fill=second!8, draw=second!50!black, left=2.4cm of mpc] (est) {状态估计\\IMU + 足底里程计 (KF)\\\scriptsize“在哪儿”};
  % control flow
  \draw[ctrl] (cmd)--(plan);
  \draw[ctrl] (plan)--(mpc) node[midway,right,font=\scriptsize]{期望轨迹};
  \draw[ctrl] (mpc)--(wbc) node[midway,right,font=\scriptsize]{$\mathbf{f}_i$（虚拟力）};
  \draw[ctrl] (wbc)--(motor) node[midway,right,font=\scriptsize]{$\boldsymbol\tau$};
  \draw[ctrl] (motor)--(robot);
  % state flow (dashed) from robot up to estimator, then to layers
  \draw[state] (robot.west) -| ([xshift=-1.2cm]robot.west) |- (est.south west);
  \draw[state] (est.north) |- ([yshift=2mm]plan.west);
  \draw[state] (est.east) -- (mpc.west);
  \draw[state] (est.south) |- ([yshift=2mm]wbc.west);
\end{tikzpicture}
\caption{四足动态运动控制的分层数据流。实线 = 控制信号流（指令逐层下传：上层指令 $\to$ 规划 $\to$ MPC $\to$ WBC $\to$ 电机 $\to$ 本体），虚线 = 状态信号流（本体状态经状态估计器反馈给规划/MPC/WBC）。三问助记标在右侧：“在哪儿”= 状态估计，“去哪儿”= 规划，“怎么去”= MPC + WBC。频率解耦：MPC 慢环 $\Delta T=0.01\,\mathrm{s}$，其余快环 $\Delta t=0.002\,\mathrm{s}$。}
\label{fig:qa-pipeline}
\end{figure}

\subsection{“在哪儿” \texorpdfstring{$=$}{=} 状态估计}
\label{subsec:qa-est}

控制器需要知道\emph{当前状态} $\boldsymbol x(0)$（位姿、速度、角速度）才能起算。它自然拆成转动与平动两部分\cite{yu2021quadruped}：
\begin{itemize}
  \item \textbf{躯干旋转}：由 \textbf{IMU} 实现——只需对 IMU 的原始数据（陀螺、加速度计）作适当的坐标变换即可得到机身姿态与角速度。
  \item \textbf{躯干平动}：用\emph{卡尔曼滤波的足底里程计}——在“支撑足不打滑”假设下，建立本体系与世界系之间的足底速度映射，从而推出质心的位置与速度。
\end{itemize}
完整的滤波器设计（ESKF / 不变滤波 RIEKF、足式里程计建模）见 \cref{ch:quad_est}，其几何与滤波基础见 \cref{ch:eskf,ch:invariant-filter,ch:imu,ch:leg}。

\subsection{“去哪儿” \texorpdfstring{$=$}{=} 规划}
\label{subsec:qa-plan}

规划层把上层速度/步态指令翻译成\emph{期望轨迹}，分两支\cite{yu2021quadruped}：
\begin{itemize}
  \item \textbf{质心轨迹规划器}：解算期望的质心位置/速度/姿态轨迹（供 MPC 跟踪）。
  \item \textbf{足底轨迹规划器}：按\emph{实际与期望水平速度的偏差}选取落脚点（落点近似为速度的线性函数，承自 Raibert，\cref{subsec:qa-milestones}），并在\emph{离地点与落脚点之间}规划摆动腿的\emph{摆线轨迹}。
\end{itemize}
完整的步态调度（接触序列、相位）、落足点选取与摆线轨迹生成见 \cref{ch:quad_gait}。

\subsection{“怎么去” \texorpdfstring{$=$}{=} MPC \texorpdfstring{$+$}{+} WBC}
\label{subsec:qa-mpc-wbc}

这是控制的核心两层，也是本部最重的两章。

\paragraph{上层 MPC：求足底反力。}
把整机\emph{简化为单刚体}（忽略腿部质量分布变化），建立状态方程与滚动预测方程，在预测时域上求解一个凸 QP，得到\emph{支撑相}各腿的最优\emph{足底反力} $\mathbf{f}_i$。单刚体建模、三处凸化近似、状态空间与 QP 化的全量推导见 \cref{ch:quad_mpc}（本部范本章）；浮动基建模与单腿运动学/雅可比见 \cref{ch:quad_kin}；QP 求解器基础见 \cref{ch:nlopt}。

\paragraph{下层 WBC：转关节力矩。}
WBC 用\emph{多刚体动力学}，把 MPC 给出的足底力/期望加速度作为\emph{任务}，按\emph{重要程度}分若干优先级（典型为接触约束 $>$ 机身位姿 $>$ 摆动腿 $>$ 其余），用\emph{雅可比零空间投影}保证高优先级任务不被低优先级扰动，最终经多刚体逆动力学解出\emph{关节力矩}发给电机。WBC 的优先级任务分解、零空间投影、浮动基逆动力学见 \cref{ch:quad_wbc}。

\begin{insight}[浮动基六自由度欠驱——躯干是被控对象却无电机直驱]\label{ins:qa-floating-base}
四足控制的\emph{结构性}难点在于：四足机器人是一个\textbf{六自由度欠驱的浮动基座}系统，而\emph{这个浮动基座（躯干）恰恰是主要的被控对象}——我们想控制的正是躯干的姿态与速度。但躯干上\emph{没有任何电机直接驱动}它的 $6$ 个自由度；唯一的执行器是\emph{关节电机}。因此，控制器只能\textbf{通过关节力矩去调节腿与环境的交互力（足底反力），再借这些力\emph{间接}维持躯干的姿态与速度}\cite{yu2021quadruped}。这就解释了为什么架构里 MPC 的决策变量是“足底力”而非“躯干力”：躯干力无法直接施加，只能由可控的足底力\emph{合成}出来。
\end{insight}

\begin{insight}[地面反力是除重力外唯一外力，故控力而非控位置]\label{ins:qa-grf-only}
紧承上一条：对一个在地面行走的浮动基机器人，\textbf{除了重力，地面反力（GRF）是它受到的\emph{唯一}外力}\cite{yu2021quadruped}。这是一个看似平凡却极其关键的事实：既然能影响躯干运动的外力\emph{只有}地面反力（重力是不可控常量），那么控制系统\emph{最该精确掌控的就是这股力}——这就从第一性原理上把范式从“控位置”定死为“\emph{控力}”（呼应 \cref{sec:qa-naive} 反面二）。MPC 之所以以足底反力为决策变量、WBC 之所以做力/力矩级的逆动力学，根子都在这里。
\end{insight}

\paragraph{足底力 $\to$ 关节力矩：最后一公里（定位式）。}
MPC/QP 在“整机/世界系”里求出了足底反力 $\mathbf{f}$，但电机只能直接控制\emph{关节力矩} $\boldsymbol\tau$。二者由\emph{单腿雅可比转置}连接——这是“足端力（世界系）$\to$ 关节空间”的最后一公里：
\begin{equation}
  \boldsymbol\tau=\mathbf{J}^\top\mathbf{f}.
  \label{eq:qa-jtf}
\end{equation}
本章\emph{只作一句话定位}，不展开虚功推导。

\begin{insight}[集各家之长——本部各章如何拼成一张图]\label{ins:qa-assemble}
把上面所有层装回 \cref{ins:qa-evolution} 的总声明，就得到本书控制系统的\emph{全貌}：\textbf{以 MPC 求足底反力（\cref{ch:quad_mpc}）+ 以 WBC 算关节空间期望并经多刚体动力学求扭矩（\cref{ch:quad_wbc}）+ 借鉴 VMC/SLIP 的力控理念}\cite{yu2021quadruped}。\cref{eq:qa-jtf} 正是 MPC 与 WBC/电机之间的接缝。由此可把后续每一章都\emph{挂}到 \cref{fig:qa-pipeline} 的某一层上：运动学/浮动基建模（\cref{ch:quad_kin}）支撑雅可比与单刚体几何；步态（\cref{ch:quad_gait}）供“去哪儿”；状态估计（\cref{ch:quad_est}）供“在哪儿”；MPC（\cref{ch:quad_mpc}）与 WBC（\cref{ch:quad_wbc}）是“怎么去”的上下两层；实践（\cref{ch:quad_prac}）把它们在宇树框架里跑通。
\end{insight}

\begin{note}[全局约定：$\mathbf{J}^\top\mathbf{f}$ 的符号与坐标系]\label{note:qa-jtf-sign}
\cref{eq:qa-jtf} 是\emph{定位性}的简写。在本书右扰动主线、机身$\to$世界旋转记 $\mathbf{R}_{sb}$ 的约定下，把\emph{世界系}足底力 $\mathbf{f}_{is}$ 转为第 $i$ 腿关节力矩的完整形式带有一个作用-反作用负号与坐标变换：
\[
  \boldsymbol\tau_i=-\,\mathbf{J}_i^\top\,\mathbf{R}_{sb}^\top\,\mathbf{f}_{is},
\]
其中负号来自作用-反作用（足端对地的力与地对足端的反力互为相反），$\mathbf{R}_{sb}^\top$ 把世界系力先转回机身系。完整的虚功推导与符号讨论见 \cref{ch:quad_prac,ch:quad_wbc}；\cref{ch:quad_mpc} 给出的逆静力学形式 $\boldsymbol\tau_i=\mathbf{J}_i^\top\mathbf{R}_i^\top\mathbf{f}_i$ 是同一关系在该章局部坐标约定下的写法。
\end{note}

\begin{note}[前向指针：这套“怎么去”分层在工业级开源栈上的实例]\label{note:qa-ocs2-pointer}
本节给出的“上层 MPC 求足底反力 $\to$ 下层 WBC 转关节力矩”分层，并非只是教学抽象：本部将把它在一个真实的工业级开源栈（OCS2 / \texttt{legged\_control}）上落成一个完整实例——\emph{质心动力学 NMPC}（保留角动量非线性，相对单刚体凸 MPC 是更高保真的一档）$+$\emph{加权 / 分层 WBC}（\cref{ch:quad_wbc} 的 WQP/HQP 谱系在代码里的两种实现）$+$\emph{线性卡尔曼滤波}（\cref{ch:quad_est} 的足式状态估计），并按本节“怎么去·上”与下文频率解耦在同一闭环里跑起来\cite{sleiman2021unified,grandia2022perceptive,liao2023walking}。这套实例的全栈数据流、移植一台自定义四足的失败经验与可迁移调试方法论，留到 \cref{ch:quad_advanced} 完整展开。
\end{note}

\subsection{频率解耦：为什么 MPC 慢、WBC 快}
\label{subsec:qa-freq}

分层的一个关键工程后果是\emph{频率解耦}：MPC 与 WBC 跑在\emph{不同}频率上\cite{yu2021quadruped}。
\begin{equation}
  \Delta T=0.01\ \text{s}\ (\text{MPC},\ 100\ \text{Hz}),
  \qquad
  \Delta t=0.002\ \text{s}\ (\text{其余模块},\ 500\ \text{Hz}).
  \label{eq:qa-freq}
\end{equation}

\begin{insight}[频率解耦——慢环看得远，快环跟得紧]\label{ins:qa-freq}
为什么不合成一层、一个频率？因为两层的\emph{计算负担}与\emph{职责}天然不同：MPC 要在线求解一个\emph{大规模、跨预测时域}的优化，算力重、只能跑\emph{慢环}（$100\,\mathrm{Hz}$），但它\emph{看得远}（为未来相位准备）；WBC/状态估计/规划只做\emph{当前周期}的逆动力学与滤波，算得轻、可以跑\emph{快环}（$500\,\mathrm{Hz}$），从而\emph{跟得紧}（高带宽抑扰、快速响应触地冲击）。慢环给出“战略”（足底力的最优分配），快环执行“战术”（高频把战略落到关节、并实时纠偏）——这正是级联/分层控制的精髓。频率与周期互为倒数：$\Delta T=0.01\,\mathrm{s}$ 对应 $100\,\mathrm{Hz}$、$\Delta t=0.002\,\mathrm{s}$ 对应 $500\,\mathrm{Hz}$（即 \cref{eq:qa-freq}）。
\end{insight}

\begin{pitfall}[误区：把所有环都拉到同一高频]
有人想“反正算力够，干脆让 MPC 也跑 $500\,\mathrm{Hz}$”。问题在于 MPC 单次求解可能就要近 $1\,\mathrm{ms}$ 量级（取决于时域与腿数），\emph{挤不进} $2\,\mathrm{ms}$ 的快环预算；强行提速只能缩短预测时域，反而\emph{丢掉了 MPC 唯一的价值——前瞻}。正确做法是\emph{分层异频}：让贵的、看得远的 MPC 慢跑，让便宜的、跟得紧的 WBC 快跑（\cref{ins:qa-freq}）。
\end{pitfall}

\begin{practice}
\begin{enumerate}
  \item 把 \cref{fig:qa-pipeline} 在你自己的四足（或仿真）上重画一遍，标注每一层的实际运行频率与输入/输出数据维度。
  \item 用 \cref{ins:qa-floating-base,ins:qa-grf-only} 回答前置自测第 $3$、$4$ 题，写成两段完整论证。
  \item （概念）若状态估计层突然漂移（“在哪儿”答错了），\cref{fig:qa-pipeline} 中哪些层会被\emph{毒化}？为什么实践中常建议“跟踪速度而非绝对位置”来削弱这一影响？
\end{enumerate}
\end{practice}

% ════════════════════════════════════════════════════════════════
\section{相位组织：有限状态机与安全机制}
\label{sec:qa-fsm}

\cref{ins:qa-multimodel} 留下的伏笔在此落地：既然四足是“多模型须切换”，就需要一个\emph{机构}来管理“此刻用哪套模型/行为、何时允许切换”。这个机构就是\textbf{有限状态机（FSM）}\cite{yu2021quadruped}。

\subsection{FSM 管理行为切换}
\label{subsec:qa-fsm-states}

FSM 把机器人的高层行为离散成若干\emph{状态}（如\emph{被动}、\emph{固定站立}、\emph{下蹲/起立}、\emph{自由站立（平衡）}、\emph{行走（trot）}等），并定义状态之间\emph{允许}与\emph{禁止}的切换。\cref{fig:qa-fsm} 给出一个典型的切换关系示意。

\begin{figure}[ht]
\centering
\begin{tikzpicture}[>=Stealth, node distance=10mm,
  st/.style={draw, rounded corners, align=center, font=\small, inner sep=3pt, minimum width=1.9cm, minimum height=0.8cm},
  tr/.style={-{Stealth[length=2mm]}, thick}]
  \node[st, fill=gray!12] (passive) {被动\\(damping)};
  \node[st, fill=main!10, right=of passive] (stand) {固定站立};
  \node[st, fill=second!10, above right=6mm and 12mm of stand] (free) {自由站立\\(平衡)};
  \node[st, fill=second!10, below right=6mm and 12mm of stand] (trot) {行走\\(trot)};
  \node[st, fill=third!12, below=14mm of passive] (safe) {阻尼保护\\(危险态)};
  % allowed transitions
  \draw[tr] (passive)--(stand);
  \draw[tr] (stand) to[bend left=10] (passive);
  \draw[tr] (stand)--(free);
  \draw[tr] (free) to[bend right=10] (stand);
  \draw[tr] (stand)--(trot);
  \draw[tr] (trot) to[bend left=10] (stand);
  \draw[tr] (free) to[bend left=10] (trot);
  \draw[tr] (trot) to[bend left=10] (free);
  % danger -> safe from any (示意：从 trot/free 触发)
  \draw[tr, third!70!black, dashed] (trot.west) to[bend right=20] node[below,font=\scriptsize,pos=0.6]{危险} (safe.east);
  \draw[tr, third!70!black, dashed] (free.west) to[bend left=25] (safe.north east);
  % forbidden (示意): squat直接 trot 用红叉
  \node[red!70!black, font=\scriptsize] at (1.7,-1.0) {$\times$ 非法直跳};
\end{tikzpicture}
\caption{四足行为有限状态机（FSM）示意。实线箭头为\emph{允许}的切换（须满足切换关系约束与“忙碌时间”），虚线箭头表示\emph{任何受控状态}在检测到危险时都可\emph{无条件}切入“阻尼保护”态。非法的直接切换（如绕过中间态从下蹲直跳行走）被禁止。}
\label{fig:qa-fsm}
\end{figure}

调度逻辑包含两条关键约束\cite{yu2021quadruped}：
\begin{itemize}
  \item \textbf{切换关系约束}：只允许在“相邻/兼容”的状态间切换。例如必须先经“固定站立”才能进入“行走”，不允许从“下蹲”\emph{直接}跳到“行走”——否则腿部构型与动力学假设不连续，极易失稳。
  \item \textbf{忙碌时间（busy time）}：某些状态切换或动作需要一段\emph{最短执行时间}才能完成（如起立动作）。在“忙碌时间”内，FSM \emph{拒绝}响应新的切换请求，等动作完成后才接受——防止动作被中途打断造成危险。
\end{itemize}

\subsection{危险检测与阻尼保护}
\label{subsec:qa-fsm-safe}

FSM 的安全落点是一个\emph{无条件优先}的“阻尼保护”态：\textbf{任何受控状态一旦检测到危险，立即被调度到阻尼保护态}\cite{yu2021quadruped}。在该态下，$12$ 个电机各自\emph{独立}施加一个与自身转速\emph{反向、且正比于转速}的力矩（纯速度阻尼），使机器人快速而平稳地静止下来：
\begin{equation}
  \boldsymbol\tau_{\mathrm{motor}}=-\,k_{\mathrm{damping}}\,v_{\mathrm{motor}}.
  \label{eq:qa-damping}
\end{equation}
其中 $k_{\mathrm{damping}}>0$ 为阻尼系数（调试确定），$v_{\mathrm{motor}}$ 为电机转速。转速越大，对外的阻尼力矩越大——这相当于给每个关节装上一个“虚拟阻尼器”，把机器人的动能平稳耗散掉，而非硬性锁死（硬锁会产生冲击）\cite{yu2021quadruped}。

\begin{insight}[阻尼保护是“多模型切换”哲学的安全落点，亦是完整关节律的退化]\label{ins:qa-damping-degenerate}
\cref{eq:qa-damping} 看似简单，却是 \cref{ins:qa-multimodel}“多模型须切换”哲学的\emph{安全落点}：当所有正常行为模型都可能失效（危险来临），系统切到一个\emph{最简单、最鲁棒}的模型——纯速度阻尼。它还是完整关节控制律的一个\emph{退化形式}。完整的电机 PD$+$前馈律一般写作
\[
  \boldsymbol\tau=k_p(\boldsymbol p_{\mathrm{des}}-\boldsymbol p)+k_d(\boldsymbol v_{\mathrm{des}}-\boldsymbol v)+\boldsymbol\tau_{\mathrm{ff}},
\]
取 $k_p=0$、$\boldsymbol v_{\mathrm{des}}=\mathbf{0}$、$\boldsymbol\tau_{\mathrm{ff}}=\mathbf{0}$、并记 $k_d=k_{\mathrm{damping}}$，即得 \cref{eq:qa-damping}。完整关节控制律及其在宇树框架中的实现见 \cref{ch:quad_prac}。
\end{insight}

\begin{table}[H]\centering\small
\caption{阻尼保护触发条件\cite{yu2021quadruped}（阈值为该平台经验值）}
\begin{tabular}{ll}
\toprule
检测量 & 触发阈值（超过即切入阻尼保护） \\
\midrule
俯仰角 + 横滚角之和（姿态倾覆） & $>35^\circ$ \\
机身水平速度 & $>5.5\ \text{m/s}$ \\
足底速度 & $>5.5\ \text{m/s}$ \\
关节速度 & $>18^\circ/\text{s}$（按平台标定） \\
关节角跟踪误差 & $>20^\circ$ \\
足底位置跟踪误差 & $>0.2\ \text{m}$ \\
\bottomrule
\end{tabular}
\label{tab:qa-danger}
\end{table}

\begin{pitfall}[误区：忽视忙碌时间或允许非法切换]
两类典型 FSM 错误都会造成危险：(i) \emph{允许非法切换}——例如让机器人从“下蹲”直接切到“行走”，腿部构型突变、动力学假设不连续，瞬间失稳；(ii) \emph{忽视忙碌时间}——在起立动作未完成时就响应“开始行走”，动作被打断、姿态未就绪即受载。正确做法是把“合法切换图 + 忙碌时间门控 + 危险检测优先级”三者一起实现（\cref{fig:qa-fsm}，\cref{tab:qa-danger}）。
\end{pitfall}

\begin{practice}
\begin{enumerate}
  \item 为你的四足（或 \texttt{unitree\_guide} 等开源框架）列一张完整的 FSM \emph{状态转移表}：行/列为状态，单元格标“允许/禁止”，并标注哪些切换有“忙碌时间”门控。
  \item 用 \cref{eq:qa-damping} 解释：为什么阻尼保护用“速度阻尼”而不是“位置锁死”来让机器人停下？后者会带来什么问题？
  \item （概念）从 \cref{tab:qa-danger} 任选两个触发条件，说明它们各自在防范什么样的失效模式（如倾覆、超速、跟踪发散）。
\end{enumerate}
\end{practice}

% ════════════════════════════════════════════════════════════════
\section{优化与强化学习：现代的两条腿}
\label{sec:qa-two-streams}

本章主体走的是\emph{优化/MPC$+$WBC}这条腿；但现代四足控制还有\emph{另一条腿}——\emph{强化学习（RL）}。本节点到为止，把二者的关系讲清，作为全章收束并指向 \cref{ch:rl-basic}。

\begin{insight}[优化与 RL 本质同源——都是在某函数下求最优运动策略]\label{ins:qa-two-streams}
有一个统一视角：\textbf{优化法与强化学习，本质上都是“在某个函数（代价/奖励）下，求最优的运动策略”}。
\begin{itemize}
  \item \emph{优化法}：显式\emph{建模}系统动力学，\emph{最小化代价函数}求未来一段时间的最优控制量（如足端力）——这就是 MPC。
  \item \emph{强化学习}：不必显式建模动力学，靠与环境/仿真器交互、\emph{基于奖励信号探索}，\emph{最大化累计回报}求一个\emph{策略}（状态到动作的映射）。
\end{itemize}
“最小化代价”与“最大化奖励”只差一个符号，二者在数学结构上是同一类问题的两种姿态\cite{yu2021quadruped}。理解这一点，就不会把“优化派”和“学习派”看成对立——它们是同一目标的两条求解路径。RL 的系统基础见 \cref{ch:rl-basic}。
\end{insight}

\begin{insight}[接触-规划分离降复杂度，但难扩展到接触不确定性 $\Rightarrow$ RL 兴起]\label{ins:qa-contact-split}
为什么 MPC 范式如此成功，却又给 RL 留出了空间？关键在 MPC 采用了\textbf{“接触与运动规划分离”}的策略：把\emph{何时何处接触}（接触序列、落足点）交给\emph{外部步态调度器}\emph{预先}给定，MPC 只在\emph{给定接触}下求最优力。这一分离\emph{大幅降低了在线优化的复杂度}（否则同时优化接触时序会变成非凸的混合整数问题），再\emph{结合 WBC}取得了很好的工程效果\cite{yu2021quadruped}。但代价是：当面对\emph{接触不确定性}（不知道脚会不会打滑、地形是否如预期）、或需要\emph{端到端融合感知}时，这种“先定接触再算力”的分离就\emph{难以扩展}。正是这一局限，催生了用 RL 端到端学习“感知$\to$动作”策略的浪潮。\textbf{MPC 的优势（分离降复杂度）与其局限（难处理接触/感知不确定性）是一枚硬币的两面}。
\end{insight}

\paragraph{模拟器是 RL 的使能器（凝练）。}
RL 之所以在近年可行，离不开\emph{大规模并行模拟器}：早期“弹簧-阻尼惩罚法”难以模拟刚性接触，仿真不准；现代 GPU 并行模拟器（如 Isaac 系列）每秒可产生\emph{数十万级}的状态转换，使 RL 所需的海量采样成为可能\cite{yu2021quadruped}。模拟器细节属工程，归入 \cref{ch:quad_prac,ch:engineering}。

\begin{table}[H]\centering\small
\caption{优化（MPC$+$WBC）与强化学习两流派对照}
\begin{tabularx}{\linewidth}{p{0.22\textwidth} L L}
\toprule
维度 & 优化 / MPC$+$WBC（本章主线） & 强化学习 RL（另一条腿） \\
\midrule
是否显式建模动力学 & 是（单刚体 / 多刚体） & 否（靠交互探索） \\
目标函数 & 最小化代价（跟踪 + 正则） & 最大化累计奖励 \\
求解物 & 当前及未来的最优控制量 & 状态到动作的策略 \\
接触处理 & 接触序列外部给定（接触-规划分离） & 可端到端学习含接触的策略 \\
主要优势 & 可解释、可加硬约束、样本零成本 & 易融合感知、处理接触不确定性 \\
主要局限 & 难扩展到接触/感知不确定性 & 需海量采样、sim-to-real 间隙 \\
代表入口 & \cref{ch:quad_mpc}、\cref{ch:quad_wbc} & \cref{ch:rl-basic} \\
\bottomrule
\end{tabularx}
\label{tab:qa-two-streams}
\end{table}

\begin{practice}
\begin{enumerate}
  \item 用 \cref{ins:qa-two-streams} 的“代价 vs 奖励差一个符号”，把一个 MPC 的跟踪代价函数改写成等价的 RL 奖励（文字描述即可）。
  \item 举一个“接触不确定性”的具体场景（如脚踩到松动石块），说明为什么 MPC 的“接触-规划分离”在此处吃力，而 RL 策略可能更鲁棒（\cref{ins:qa-contact-split}）。
  \item （开放）查阅一个把 MPC 与 RL \emph{混合}使用的工作，说明它分别让两条腿负责什么。
\end{enumerate}
\end{practice}

% ════════════════════════════════════════════════════════════════
\section{章节定位总表：后续每一章在架构图上的位置}
\label{sec:qa-map}

作为全章的收束与“地图索引”，\cref{tab:qa-chapter-map} 把本部后续每一章\emph{钉}在 \cref{fig:qa-pipeline} 的对应层上，并标出它回答三问中的哪一问。读任意一章时回看此表，即知“我在地图的哪个位置”。

\begin{table}[H]\centering\small
\caption{本部各章在分层架构中的定位（\cref{fig:qa-pipeline} 索引）}
\begin{tabular}{p{0.30\textwidth} p{0.16\textwidth} p{0.42\textwidth}}
\toprule
章节 & 三问归属 & 在架构中的角色 \\
\midrule
\cref{ch:quad_arch}（本章）导论与架构 & 全局 & 地图本身：建立分层数据流与控制哲学 \\
\cref{ch:quad_kin} 浮动基建模与运动学 & 基础 & 提供浮动基模型、单腿雅可比 $\mathbf{J}$（供 \cref{eq:qa-jtf}）\\
\cref{ch:quad_gait} 步态生成与轨迹规划 & 去哪儿 & 规划层：接触序列、落足点、摆线轨迹 \\
\cref{ch:quad_est} 足式状态估计 & 在哪儿 & 状态估计层：IMU + 足底里程计给出 $\boldsymbol x(0)$ \\
\cref{ch:quad_mpc} 单刚体 MPC & 怎么去·上 & MPC 慢环：单刚体 $\to$ 凸 QP $\to$ 足底反力 $\mathbf{f}_i$ \\
\cref{ch:quad_wbc} 全身控制 WBC & 怎么去·下 & WBC 快环：多刚体逆动力学 $\to$ 关节力矩 $\boldsymbol\tau$ \\
\cref{ch:quad_prac} 实践：宇树框架 & 落地 & FSM、电机控制律、整定、部署（\cref{ch:engineering}）\\
\bottomrule
\end{tabular}
\label{tab:qa-chapter-map}
\end{table}

% ════════════════════════════════════════════════════════════════
\section*{本章常见误解汇总}

\begin{table}[H]\centering\small
\begin{tabular}{p{0.46\textwidth} p{0.46\textwidth}}
\toprule
\textbf{常见误解} & \textbf{正确理解} \\
\midrule
四足控制和汽车控制只是“腿多了几条” & 汽车是连续单模型，四足是多模型须切换（\cref{ins:qa-multimodel}）\\
控制就是设计 PID 把误差压下去 & 控制本质是带约束优化；难解时换次优可行解（\cref{ins:qa-control-is-opt}）\\
位置控制足以驱动四足 & 与环境有力交互时位置环是黑箱、不可控（\cref{sec:qa-naive} 反面二）\\
纯 WBC/瞬时 QP 足以做动态步态 & 无预测性，相位切换处无法准备，须 MPC（\cref{ins:qa-need-mpc}）\\
躯干由专门的电机驱动 & 浮动基六自由度欠驱，躯干无直驱，靠足底力间接维持（\cref{ins:qa-floating-base}）\\
该控躯干的位置 & 除重力外地面反力是唯一外力，故控力（\cref{ins:qa-grf-only}）\\
MPC 与 WBC 应合成一层同频率 & 慢环看得远、快环跟得紧，频率解耦（\cref{ins:qa-freq}）\\
MPC 直接输出关节力矩 & MPC 出足底虚拟力，再 $\boldsymbol\tau=\mathbf{J}^\top\mathbf{f}$ 转力矩（\cref{eq:qa-jtf}）\\
一套模型 + 一个控制器走天下 & 须 FSM 管理多行为合法切换（\cref{sec:qa-fsm}）\\
危险时把关节锁死最安全 & 锁死有冲击；用速度阻尼平稳耗能（\cref{eq:qa-damping}）\\
优化派与学习派是对立路线 & 都是在某函数下求最优策略，同源（\cref{ins:qa-two-streams}）\\
SLIP 公式是四足的精确动力学 & \cref{eq:qa-slip} 是点质量概念模型，仅供直觉（\cref{rem:qa-slip}）\\
\bottomrule
\end{tabular}
\end{table}

% ════════════════════════════════════════════════════════════════
\section*{本章小结}

\begin{table}[H]\centering\small
\caption{符号 / 概念表}
\begin{tabular}{lll}
\toprule
符号 / 概念 & 含义 & 首次出现 \\
\midrule
控制 $=$ 优化/规划 & 目标 + 约束 + 变量；难解换次优可行 & \cref{ins:qa-control-is-opt} \\
连续单模型 vs 多模型 & 汽车/无人机 vs 四足须切换 & \cref{ins:qa-multimodel} \\
三问 & 在哪儿/去哪儿/怎么去 = 估计/规划/控制 & \cref{ins:qa-three-questions} \\
浮动基欠驱 & 六自由度躯干无直驱，靠足底力间接控 & \cref{ins:qa-floating-base} \\
GRF $\mathbf{f}_i$ & 第 $i$ 腿足底反力（除重力外唯一外力） & \cref{ins:qa-grf-only} \\
$\boldsymbol\tau=\mathbf{J}^\top\mathbf{f}$ & 足底力 $\to$ 关节力矩（最后一公里） & \cref{eq:qa-jtf} \\
$\mathbf{R}_{sb}$ & 机身$\to$世界旋转 & \cref{note:qa-jtf-sign} \\
$\Delta T,\Delta t$ & MPC 慢环 / 其余快环周期 & \cref{eq:qa-freq} \\
FSM & 有限状态机，管理行为合法切换 & \cref{sec:qa-fsm} \\
$\boldsymbol\tau_{\mathrm{motor}}=-k_{\mathrm{damping}}v_{\mathrm{motor}}$ & 阻尼保护控制律 & \cref{eq:qa-damping} \\
SLIP（示意） & 弹簧负载倒立摆概念模型 & \cref{eq:qa-slip} \\
优化 vs RL & 最小化代价 vs 最大化奖励，同源 & \cref{ins:qa-two-streams} \\
\bottomrule
\end{tabular}
\end{table}

\begin{table}[H]\centering\small
\caption{要点速查表}
\begin{tabular}{lll}
\toprule
要点 / 结论 & 一句话说明 & 对应 \\
\midrule
控制即优化 & 目标 + 约束 + 变量；牺牲解空间换可行 & \cref{ins:qa-control-is-opt} \\
多模型须切换 & 四足非连续单模型 $\Rightarrow$ FSM & \cref{ins:qa-multimodel,sec:qa-fsm} \\
四反面 & 连杆/位置/瞬时/早期学习各自的失败 & \cref{sec:qa-naive} \\
六里程碑 & 连杆$\to$位置$\to$人驾$\to$Raibert$\to$VMC$\to$WBC/MPC & \cref{fig:qa-timeline} \\
三问助记 & 在哪儿/去哪儿/怎么去 & \cref{ins:qa-three-questions} \\
分层数据流 & 估计$\to$规划$\to$MPC$\to$WBC$\to$电机 & \cref{fig:qa-pipeline} \\
欠驱 $\to$ 力控 & 躯干无直驱、GRF 唯一外力 & \cref{ins:qa-floating-base,ins:qa-grf-only} \\
$\boldsymbol\tau=\mathbf{J}^\top\mathbf{f}$ & 足底力转关节力矩 & \cref{eq:qa-jtf} \\
频率解耦 & MPC $100\,\mathrm{Hz}$ / 其余 $500\,\mathrm{Hz}$ & \cref{eq:qa-freq} \\
阻尼保护 & 危险态纯速度阻尼平稳停机 & \cref{eq:qa-damping} \\
两流派同源 & 优化与 RL 都是求最优策略 & \cref{ins:qa-two-streams} \\
接触-规划分离 & 降复杂度但难扩展接触不确定性 & \cref{ins:qa-contact-split} \\
\bottomrule
\end{tabular}
\end{table}

\begin{table}[H]\centering\small
\caption{知识点总表}
\begin{tabular}{clll}
\toprule
\# & 知识点 & 核心要点 & 难度 \\
\midrule
1 & 控制即优化/规划 & 目标 + 约束 + 变量；难解换次优可行 & \(\bigstar\) \\
2 & 连续单模型 vs 多模型 & 四足须切换 $\Rightarrow$ FSM 必要性 & \(\bigstar\bigstar\) \\
3 & 朴素思路四反面 & 逼出力控/预测/分层/学习 & \(\bigstar\bigstar\) \\
4 & 六里程碑控制史 & 从无反馈到 MPC+WBC 的演进主轴 & \(\bigstar\bigstar\) \\
5 & 典型平台谱 & BD/MIT/ETH/IIT/UPenn/宇树 & \(\bigstar\) \\
6 & 三问与分层架构 & 估计/规划/MPC/WBC/电机数据流 & \(\bigstar\bigstar\) \\
7 & 浮动基欠驱与力控 & 躯干无直驱、GRF 唯一外力 & \(\bigstar\bigstar\bigstar\) \\
8 & 足底力转关节力矩 & $\boldsymbol\tau=\mathbf{J}^\top\mathbf{f}$ 定位式 & \(\bigstar\bigstar\) \\
9 & 频率解耦 & 慢环看得远、快环跟得紧 & \(\bigstar\bigstar\) \\
10 & FSM 与阻尼保护 & 合法切换 + 忙碌时间 + 危险阻尼 & \(\bigstar\bigstar\) \\
11 & 优化 vs RL 两流派 & 同源；接触-规划分离的得与失 & \(\bigstar\bigstar\) \\
\bottomrule
\end{tabular}
\end{table}

% ════════════════════════════════════════════════════════════════
\section*{延伸阅读}
\begin{itemize}
  \item \textbf{一手分层框架与 FSM}：于宪元《基于稳定性的仿生四足机器人控制系统设计》\cite{yu2021quadruped}（北航硕论，2021）——本章控制史六里程碑、浮动基欠驱本质、三问助记、MPC$+$WBC 分层与 FSM 实现的直接出处。
  \item \textbf{工程落地与平台谱}：宇树《四足机器人控制算法》\cite{unitree2023quadruped}——典型平台谱、四系统组成、$12$ 关节永磁同步电机硬件前提，配套 \texttt{unitree\_guide} 开源代码（\cref{ch:quad_prac}）。
  \item \textbf{奠基范式}：Raibert\cite{raibert1986legged}（解耦三通道 / SLIP，证明“精准动力学非必要”）、Pratt 等\cite{pratt2001virtual}（虚拟模型控制 VMC）——历史里程碑 \cnum{4}、\cnum{5} 的原始文献。
  \item \textbf{现代凸 MPC 范式}：Di Carlo 等\cite{dicarlo2018cheetah}（IROS 2018，单刚体凸 MPC，硬件平台 Cheetah 3\cite{bledt2018cheetah3}）——“怎么去·上”那一层的标准范式，本部范本章 \cref{ch:quad_mpc} 的主线。
  \item \textbf{简化模型谱系与综述}：Kajita\cite{kajita2001lipm}（LIPM）、Orin\cite{orin2013centroidal}（质心动力学）、Wensing 等\cite{wensing2024legged}（T-RO 2024 综述）——理解“架构里 MPC 那一层为何能用简化模型”的谱系背景。
\end{itemize}

\section*{本章与相邻章节的关系}
\begin{table}[H]\centering\small
\begin{tabular}{lll}
\toprule
相关章节 & 关系 & 本章接口 \\
\midrule
\cref{ch:control} 控制导论 & 提供“代价 + 约束 + 滚动优化”一般框架 & \cref{ins:qa-control-is-opt} \\
\cref{ch:nlopt} 非线性优化/QP & MPC/WBC 落地为可解问题的求解器基础 & \cref{subsec:qa-mpc-wbc} \\
\cref{ch:quad_mpc} 单刚体 MPC & “怎么去·上”层，足底反力 $\mathbf{f}_i$ & \cref{fig:qa-pipeline,ins:qa-need-mpc} \\
\cref{ch:quad_wbc} 全身控制 WBC & “怎么去·下”层，关节力矩 $\boldsymbol\tau$ & \cref{eq:qa-jtf,note:qa-jtf-sign} \\
\cref{ch:quad_est} 足式状态估计 & “在哪儿”层，给出 $\boldsymbol x(0)$ & \cref{subsec:qa-est} \\
\cref{ch:quad_gait} 步态生成与规划 & “去哪儿”层，落足点与摆线 & \cref{subsec:qa-plan} \\
\cref{ch:rl-basic} 强化学习导论 & 现代“另一条腿” & \cref{ins:qa-two-streams} \\
\bottomrule
\end{tabular}
\end{table}
